\section*{Abstract}
\addcontentsline{toc}{section}{Abstract}

Semantic communication shifts the objective of communication systems from bit-level fidelity to meaning-level fidelity, making it a promising direction for post-Shannon wireless communication. This project investigates a DeepSC-style Transformer-based semantic transceiver for text transmission over AWGN and Rayleigh fading channels. The study focuses on three improvement directions: SNR-aware attention, adaptive semantic-loss weighting, and language-model-based post-refinement. A baseline system is first reproduced, after which several modified variants are evaluated under the same data split, SNR grid, and metric definitions.

The results show that the improved transceiver gives measurable gains under low-SNR AWGN conditions. At \SI{-10}{dB}, BLEU-4 increases from 0.4189 to 0.4329, and Word Accuracy increases from 0.4425 to 0.4773. However, the semantic-loss ablation shows that removing the semantic-loss term improves BLEU-4 across all tested SNR points, indicating that the current semantic regularization is not well aligned with token-level reconstruction. Under Rayleigh fading, the gains are limited and mixed, suggesting that scalar SNR conditioning is insufficient for strongly time-varying channels. Confidence-driven BERT refinement also fails to provide reliable net gains because of over-correction. In comparison, beam search decoding gives the most consistent inference-stage improvement, reaching BLEU-4 = 0.4552 at \SI{-10}{dB} without retraining.

Overall, the project suggests that robust channel-aware design and decoding-strategy optimization are more reliable than unconstrained external post-refinement in this setting.

\vspace{0.5em}
\noindent\textbf{Keywords:} Semantic Communication; DeepSC; SNR-Adaptive Attention;
Adaptive Semantic Loss; BERT Masked Refinement; Beam Search; Rayleigh Fading;
Transformer.